\documentclass[10pt,a4paper]{amsart}
\usepackage[margin=19mm]{geometry}
\usepackage{amsmath,amssymb,mathtools}
\usepackage[scr=rsfs,cal=euler]{mathalfa}
\usepackage{xfrac}
\usepackage[unicode,hidelinks,bookmarks=false]{hyperref}
\newcommand{\ie}{i.e.\ }
\newcommand{\cf}{cf.\ }
\newcommand{\resp}{respectively\ }
\newcommand{\Subsection}{\subsection}
\providecommand{\nobreakdash}{\nobreak\hbox{-}}
\setlength{\emergencystretch}{2em}
\multlinegap0pt
\usepackage{bm}
\usepackage{bbm}
\usepackage{dsfont}
\usepackage{xspace}
\usepackage{cancel}
\usepackage{mathtools}
\usepackage[shortlabels]{enumitem}
\usepackage{amsthm}
\makeatletter
\@ifundefined{theorem}{\newtheorem{theorem}{Theorem}}{}
\@ifundefined{lemma}{\newtheorem{lemma}[theorem]{Lemma}}{}
\makeatother
\makeatletter
\newcommand{\HS}{\mathrm{HS}}
\expandafter\ifx\csname oldproof\endcsname\relax
\newenvironment{oldproof}{\par\smallskip{\color{gray}\textbf{Old proof (kept for reference).}}\ \color{gray}}{\par\smallskip}
\fi
\expandafter\ifx\csname oldstatement\endcsname\relax
\newenvironment{oldstatement}{\par\smallskip\noindent{\color{gray}\textbf{Old statement (kept for reference).}}\ \color{gray}}{\par\smallskip}
\fi
\makeatother
\title[Beurling--Kato theory, Hardy--Sobolev calculus and Ritt oper]{Beurling--Kato theory, Hardy--Sobolev calculus and Ritt operators}
\author{A. Borichev, A. Gomilko,, Yu. Tomilov}
\date{Source: arXiv:2607.10507v2 (CC BY 4.0)}
\begin{document}
\maketitle

\noindent\emph{Source and licence.} Beurling--Kato theory, Hardy--Sobolev calculus and Ritt operators. Version arXiv:2607.10507v2 (CC BY 4.0), distributed under the Creative Commons Attribution 4.0 International licence, as recorded in the arXiv licence metadata for this exact version and shown on the abstract page https://arxiv.org/abs/2607.10507v2. Full rights notice: licenses/arxiv-2607.10507-CC-BY-4.0.txt.\par\smallskip

\noindent\textit{Editorial scope.} Counted unit: the Lemma whose source label is \texttt{lem:rational-truncation-compatibility} in the article (source0.tex lines 2773--2795, source label \texttt{lem:rational-truncation-compatibility}), delivered together with the article's own complete original proof of it (source0.tex lines 2796--3196, one \texttt{proof} environment of 401 lines). No mathematics is written, repaired or re-derived: statement and proof are the article's own text line for line. Required context retained uncounted: lem:geometry (lines 1469--1545); curve (lines 1519--1523); lem:Fmu-uniform (lines 2210--2220); thm:scalar-reproducing (lines 2262--2271). Every definition, display and label of those lines is retained unchanged. Omitted from this package and not redistributed: 1--1468, 1524--2209, 2221--2261, 2272--2772, 3197--6584. Nothing inside a retained excerpt is omitted or altered: every retained body is delivered exactly as the article's lines are written, so no omission receipt of any kind accompanies this package. Citations: the retained excerpts carry no citation command at all, so no bibliography entry of the article is transcribed and no reference is redistributed. Reference adaptation: none was needed and none was made; every reference inside the delivered unit resolves inside it, so no unresolved reference or citation is produced. Printed slips kept literal: no printed word, symbol, hypothesis or number of the counted statement or of its proof was altered. Layout and class: the article's own class is \texttt{amsart} with packages including amsmath, amscd, amssymb, eucal, enumerate, color, tikz, xcolor, enumitem, comment, bm; the wrapper is the lane's pinned \texttt{amsart} editorial wrapper at 10pt on A4, which restates the theorem-like environments the retained text needs and adds the article's own macro definitions that the retained text needs (\texttt{\textbackslash HS}), with extra wrapper packages (bm, bbm, dsfont, xspace, cancel, mathtools, enumitem). The delivered numbering is the unit's own. Assets: no retained span contains a figure environment, a graphics-inclusion command, a PDF-page inclusion, a table or a TikZ picture; no image file is copied into this package, the package declares no asset dependency and no image file is redistributed with it. Licence: version arXiv:2607.10507v2 of this article is distributed under the Creative Commons Attribution 4.0 International licence, as recorded in the arXiv licence metadata for this exact version and shown on the abstract page https://arxiv.org/abs/2607.10507v2. A full rights notice is delivered at licenses/arxiv-2607.10507-CC-BY-4.0.txt. Characters: every retained excerpt is pure ASCII, so the delivered engine, XeLaTeX with its default Latin Modern text font, reports no missing character.
\begin{lemma}\label{lem:geometry}
Let \(\alpha>1\).
\begin{enumerate}[label=\textnormal{(\roman*)}]
\item If \(1-\mu=r e^{i\theta}\in\partial\Omega_\alpha\), with
\(r>0\) and \(\theta\in(-\pi/2,\pi/2)\), then
\begin{equation}\label{eq:costheta}
        \cos\theta
        =
        \frac1\alpha+\frac r2\Bigl(1-\frac1{\alpha^2}\Bigr).
\end{equation}
In particular \(\theta\to\pm\psi_\alpha\) as \(r\downarrow0\).
\item
\[
        \Omega_\alpha\subset B(0,R_\alpha)\cap \Sigma_{\psi_\alpha}.
\]
For \(0<r<R_\alpha\), put
\[
        \Theta_\alpha(r)
        :=\{\vartheta\in(-\pi/2,\pi/2):\ r e^{i\vartheta}\in\Omega_\alpha\}.
\]
Then \(\Theta_\alpha(r)=(-\theta_\alpha(r),\theta_\alpha(r))\), where
\(\theta_\alpha(r)\in[0,\psi_\alpha)\) is determined by
\[
        \cos\theta_\alpha(r)
        =
        \frac1\alpha+\frac r2\Bigl(1-\frac1{\alpha^2}\Bigr)
\]
whenever the right-hand side is at most \(1\), and
\(\Theta_\alpha(r)=\varnothing\) otherwise.
\item 
%There exists \(C>0\), independent of
%\(\alpha\), such that
One has
\begin{equation}\label{eq:Stolz-local-length-main}
        \operatorname{length}\bigl(\Gamma_\alpha\cap\{|z-1|<t\}\bigr)
        \le 4 \pi C_{\rm B}t,\qquad t>0.
\end{equation}
Furthermore, for \(0<s<t\le 1\),
\[
        \int_{\Gamma_\alpha\cap\{s\le |z-1|<t\}}
             \frac{|dz|}{|z-1|}
        \le 2\sqrt2 \log \frac{t}{s}.
\]
%\item
\item Let \(1<\tau<\alpha\). For \(\eta>0\), set
\[
        A_\eta:=\partial D(1,\eta)\setminus S_\alpha.
\]
Then, for all sufficiently small \(\eta>0\), \(A_\eta\) is a closed arc
of length at most \(2\pi\eta\), and
\begin{equation}\label{curve}
        \Gamma_{\alpha,\eta}
        :=
        \bigl(\Gamma_\alpha\cap\{|z-1|\ge\eta\}\bigr)\cup A_\eta
\end{equation}
is a rectifiable Jordan curve
whose bounded complementary component
contains \(\overline{S_\tau}\).
%the bounded complementary component of
%surrounding \(\overline{S_\tau}\).

 %Let \(1<\tau<\alpha\).  For all sufficiently small \(\eta>0\), let
%\(a_\eta^\pm\) be the two points of
%\(\Gamma_\alpha\cap\partial D(1,\eta)\), and let \(A_\eta\) be the component of
%\(\partial D(1,\eta)\setminus S_\alpha\) joining them.  Then
%\(\operatorname{length}(A_\eta)\le2\pi\eta\), and
%\begin{equation}\label{curve}
 %       \Gamma_{\alpha,\eta}:=
 %       \bigl(\Gamma_\alpha\cap\{|z-1|\ge\eta\}\bigr)\cup A_\eta
%\end{equation}
%is a rectifiable Jordan curve and satisfies
%\[
 %       \operatorname{Ind}_{\Gamma_{\alpha,\eta}}(z)=1,\qquad
 %       z\in\overline{S_\tau}.
%\]
\end{enumerate}
\end{lemma}

\begin{lemma}\label{lem:Fmu-uniform}
Let \(1<\delta<\tau\).  There exists %a constant
\(C_{\delta,\tau}>0\) such that
\begin{equation}\label{Fmu}
        \sup_{\mu\in\Gamma_\tau\setminus\{1\}}
        \sup_{z\in S_\delta}
        |L_\mu(z)|
        \le
        C_{\delta,\tau}.
\end{equation}
\end{lemma}

\begin{theorem}\label{thm:scalar-reproducing}
Let \(1<\delta<\tau\), and let \(f\in\HS(S_\tau)\).  Then, for every
\(z\in S_\delta\),
\begin{equation}\label{eq:scalar-reproducing}
        f(z)-f(1)
        =
        \frac{1}{2\pi i}
        \int_{\Gamma_\tau}f'(\mu)L_\mu(z)\,d\mu.
\end{equation}
\end{theorem}

\begin{lemma}\label{lem:rational-truncation-compatibility}
Let $1<\sigma<\tau$, and let $T$ be a Ritt operator of Stolz type
$\sigma$.  Let $f$ be holomorphic in a neighbourhood of $\overline{S_\tau}$.
For $\varepsilon>0$ put
\[
        K_\varepsilon
        :=
        \{\mu\in\Gamma_\tau:\ |\mu-1|\ge\varepsilon\},
\]
with the inherited orientation, and define
\[
        h_\varepsilon(z)
        :=
        \frac1{2\pi i}
        \int_{K_\varepsilon}f'(\mu)L_\mu(z)\,d\mu .
\]
Then
\[
        h_\varepsilon(T)\longrightarrow f(T)-f(1),
        \qquad \varepsilon\downarrow0,
\]
in the uniform operator topology, where $f(T)$ is defined by the Riesz--Dunford functional calculus.
\end{lemma}

\begin{proof}
Choose \(\delta_0\) and \(\delta\) such that
\[
        \sigma<\delta_0<\delta<\tau .
\]
All constants below may depend on \(T,\delta_0,\delta,\tau\) and \(f\), but
not on \(\varepsilon\), and they may vary from line to line.  Put
\begin{equation}\label{eta}
        \eta=\eta(\varepsilon):=\varepsilon^2 .
\end{equation}
For \(\varepsilon>0\) sufficiently small, let
\begin{equation}\label{gammacur}
%\(
        \mathcal C_\varepsilon:=\Gamma_{\delta_0,\eta(\varepsilon)}
%        \)
\end{equation}
be the positively oriented truncated curve defined as in \eqref{curve}.
Its bounded complementary component contains
\(\overline{S_\sigma}\). Moreover, 
\(\mathcal C_\varepsilon \subset \rho(T)\) and satisfies
%\[
 %       \|(\lambda -T)^{-1}\|
 %       \le \frac {C}{|\lambda-1|},
 %       \qquad \lambda\in\Gamma_\varepsilon.       
%\]
\[
        \|(\lambda-T)^{-1}\|
        \le
        \frac{M_{\delta_0}(T)}{|\lambda-1|},
        \qquad \lambda\in\mathcal C_\varepsilon.
\]
For \(\varepsilon\) small, the compact set bounded by
\(\mathcal C_\varepsilon\) is contained in the domain of holomorphy of \(f\).


%We first note that \(h_\varepsilon\) is holomorphic in a neighbourhood of
%this compact set.  Indeed, if \(\mu\in K_\varepsilon\), then the cut
%\(\gamma_\mu\) satisfies
%\[
 %       |\xi-1|\ge |\mu-1|\ge\varepsilon,
 %       \qquad \xi\in\gamma_\mu ,
%\]
%whereas the closing part of \(\mathcal C_\varepsilon\) lies on
%\[
 %       |\lambda-1|=\eta=\varepsilon^2 .
%\]
%The remaining part of \(\mathcal C_\varepsilon\) lies on
%\(\Gamma_{\delta_0}\), while the cuts \(\gamma_\mu\) lie outside \(S_\tau\).
%Thus the cuts do not meet the compact set bounded by
%\(\mathcal C_\varepsilon\).  

We first note that \(h_\varepsilon\) is holomorphic in a neighbourhood
of the compact set enclosed by \(\mathcal C_\varepsilon\).
Indeed, this compact set is contained in
\(
        \overline{S_{\delta_0}}\cup\overline{D(1,\eta)}.
\)
For \(\mu\in K_\varepsilon\), the cut \(\gamma_\mu\) lies outside
\(S_\tau\) and satisfies
\[
        |\xi-1|\ge|\mu-1|\ge\varepsilon>\eta,
        \qquad \xi\in\gamma_\mu.
\]
Since
\[
        \overline{S_{\delta_0}}\setminus\{1\}\subset S_\tau,
\]
the cuts do not meet the compact set enclosed by
\(\mathcal C_\varepsilon\).
Hence the Riesz--Dunford functional calculus gives
\[
        h_\varepsilon(T)-f(T)+f(1)
        =
        \frac{1}{2\pi i}
        \int_{\mathcal C_\varepsilon}
        E_\varepsilon(\lambda)(\lambda -T)^{-1}\,d\lambda ,
\]
where
\[
        E_\varepsilon(\lambda)
        :=
        h_\varepsilon(\lambda)-\bigl(f(\lambda)-f(1)\bigr).
\]

%On \(S_\delta\), the scalar reproducing formula gives
%\[
%        E_\varepsilon(\lambda)
%        =
 %       -\frac1{2\pi i}
 %       \int_{\Gamma_\tau\setminus K_\varepsilon}
%        f'(\mu)L_\mu(\lambda)\,d\mu.
%\]

By Theorem~\ref{thm:scalar-reproducing}, for every
$\lambda\in S_\delta$,
\[
\begin{aligned}
    E_\varepsilon(\lambda)
    &=
    h_\varepsilon(\lambda)
    -
    \bigl(f(\lambda)-f(1)\bigr)
    \\
    &=
    -\frac{1}{2\pi i}
    \int_{\Gamma_\tau\setminus K_\varepsilon}
         f'(\mu)L_\mu(\lambda)\,d\mu.
\end{aligned}
\]
%Since \(f'\) is bounded near \(1\), Lemma~\ref{lem:geometry}\textup{(iii)}
%and Lemma~\ref{lem:Fmu-uniform} imply
%\[
 %       \sup_{\lambda\in S_\delta}|E_\varepsilon(\lambda)|
 %       \le
 %       C
 %       \int_{\Gamma_\tau\cap\{|\mu-1|<\varepsilon\}} |d\mu|
 %       \le C\varepsilon .
%\]

%Hence, 
%putting
%\[
 %       M_f:=\sup_{\mu\in\Gamma_\tau}|f'(\mu)|,
%\]
%Lemma~\ref{lem:geometry}\textup{(iii)} and
%Lemma~\ref{lem:Fmu-uniform} imply
%\[
%\begin{aligned}
  %      \sup_{\lambda\in S_\delta}|E_\varepsilon(\lambda)|
   %     \le
   %     \frac{M_fC_{\delta,\tau}}{2\pi}
  %      \operatorname{length}
  %      \bigl(\Gamma_\tau\cap\{|\mu-1|<\varepsilon\}\bigr) 
 %       \le
 %       2C_{\rm B}C_{\delta,\tau}M_f\,\varepsilon .
%\end{aligned}
%\]

Put
\[
    M_f:=\sup_{\mu\in\Gamma_\tau}|f'(\mu)|.
\]
Applying Lemma~\ref{lem:Fmu-uniform} to the preceding integral, and then
using Lemma~\ref{lem:geometry} \textup{(iii)}, we obtain
\[
\begin{aligned}
    \sup_{\lambda\in S_\delta}|E_\varepsilon(\lambda)|
    &\le
    \frac{M_fC_{\delta,\tau}}{2\pi}
    \operatorname{length}
    \bigl(
        \Gamma_\tau\cap\{|\mu-1|<\varepsilon\}
    \bigr)
    \\
    &\le
    2C_{\rm B}C_{\delta,\tau}M_f\,\varepsilon.
\end{aligned}
\]


For \(\eta=\eta(\varepsilon)\) as in \eqref{eta}, write
\[
   \mathcal C_\varepsilon
   =
   \bigl(\Gamma_{\delta_0}\cap\{|\lambda-1|\ge\eta\}\bigr)
   \cup A_\eta .
\]
The first part of the union is contained in \(S_\delta\).  Therefore, by Lemma~\ref{lem:geometry}, 
\begin{gather}\label{delta0}
%\begin{aligned}
   \int_{\Gamma_{\delta_0}\cap\{|\lambda-1|\ge\eta\}}
   |E_\varepsilon(\lambda)|
   \|(\lambda-T)^{-1}\|\,|d\lambda| 
    \le
   C\varepsilon
   \int_{\Gamma_{\delta_0}\cap\{|\lambda-1|\ge\eta\}}
   \frac{|d\lambda|}{|\lambda-1|} \\
   \le
   C\varepsilon
   \Bigl(
      2\sqrt2\log\frac1\eta
      +
      \operatorname{length}(\Gamma_{\delta_0})
   \Bigr) 
   =
   C\varepsilon
   \Bigl(
      4\sqrt2\log\frac1\varepsilon
      +
      \operatorname{length}(\Gamma_{\delta_0})
   \Bigr).\notag
%\end{aligned}
\end{gather}
Thus the contribution of this part tends to zero as $\varepsilon \downarrow 0.$


%\[
 %       B_\varepsilon
 %       :=
 %       \Gamma_{\delta_0}\cap\{|\lambda-1|\ge\eta\},
  %      \qquad
  %      A_\eta
 %       :=
 %       \Gamma_\varepsilon\setminus B_\varepsilon .
%\]
%Then \(B_\varepsilon\subset S_\delta\).  Therefore
%\[
%\begin{aligned}
 %       \int_{B_\varepsilon}
 %       |E_\varepsilon(\lambda)|
 %       \|(\lambda -T)^{-1}\|\,|d\lambda|
 %       &\le
 %       C\varepsilon
 %       \int_{\Gamma_{\delta_0}\cap\{|\lambda-1|\ge\eta\}}
  %      \frac{|d\lambda|}{|\lambda-1|}       \\
 %       &\le
  %      C\varepsilon\log \frac1{\eta}
 %       \le
 %       C\varepsilon\log \frac1{\varepsilon}.
%\end{aligned}
%\]
%Thus the contribution of \(B_\varepsilon\) tends to zero.

It remains to estimate the contribution of the closing arc \(A_\eta\).  Since \(L_\mu(1)=i\pi\),
\[
        h_\varepsilon(1)
        =
        \frac12 
        \int_{K_\varepsilon}f'(\mu)\,d\mu .
\]
As \(f'\) is holomorphic in a neighbourhood of \(\overline{S_\tau}\),
\[
        \int_{\Gamma_\tau}f'(\mu)\,d\mu=0,
\]
hence
%\[
 %       |h_\varepsilon(1)|
 %       \le
 %       C
 %       \int_{\Gamma_\tau\cap\{|\mu-1|<\varepsilon\}} |d\mu|
 %       \le C\varepsilon .
%\]
\[
\begin{aligned}
        |h_\varepsilon(1)|
        \le
        \frac{M_f}{2}
        \operatorname{length}
        \bigl(\Gamma_\tau\cap\{|\mu-1|<\varepsilon\}\bigr) 
        \le
        2\pi C_{\rm B}M_f\,\varepsilon.
\end{aligned}
\]


We next estimate \(h_\varepsilon'\) near \(1\).  If
\(|\zeta-1|\le\eta\) and \(\mu\in K_\varepsilon\), then, for
\(\varepsilon\) small,
\[
        |\zeta-1|\le\eta=\varepsilon^2
        \le \frac12|\mu-1|.
\]
%Consequently,
%\[
 %       |\mu-\zeta|\ge \frac12|\mu-1|,
 %       \qquad
 %       |1-\mu\zeta|\ge \frac12|\mu-1|.
%\]
Consequently,
\begin{equation}\label{l1}
        |\mu-\zeta|
        \ge
        |\mu-1|-|\zeta-1|
        \ge
        \frac12|\mu-1|,
\end{equation}
and, since \(|\mu|\le1\),
\begin{gather}\label{l2}
%\begin{aligned}
        |1-\mu\zeta|
        =
        |(1-\mu)+\mu(1-\zeta)|  \\
        \ge
        |\mu-1|-|\mu|\,|\zeta-1|
        \ge
        \frac12|\mu-1|.\notag
%\end{aligned}
\end{gather}


Using that
\[
        \partial_\zeta L_\mu(\zeta)
        =
        \frac{1-\mu^2}{(1-\mu\zeta)(\mu-\zeta)}
\]
and
\(
        |1-\mu^2|
        \le
        2|\mu-1|, 
\)
along with the two lower bounds \eqref{l1} and \eqref{l2} proved above, we obtain
\[
        |\partial_\zeta L_\mu(\zeta)|
        \le
        \frac{8}{|\mu-1|},
        \qquad
        |\zeta-1|\le\eta,\quad \mu\in K_\varepsilon.
\]
Therefore, arguing as in \eqref{delta0},
\[
\begin{aligned}
   \sup_{|\zeta-1|\le\eta}|h_\varepsilon'(\zeta)|
   &\le
   \frac4\pi
   \int_{K_\varepsilon}
        \frac{|f'(\mu)|}{|\mu-1|}\,|d\mu| \\
   &\le
   \frac{4M_f}{\pi}
   \left(
      2\sqrt2\log\frac1\varepsilon
      +
      \operatorname{length}(\Gamma_\tau)
   \right).
\end{aligned}
\]
%\[
%\begin{aligned}
%   \sup_{|\zeta-1|\le\eta}|h_\varepsilon'(\zeta)|
 %  &\le
 %  C\int_{K_\varepsilon}
 %       \frac{|f'(\mu)|}{|\mu-1|}\,|d\mu| \\
 %  &\le
 %  C\log\left(\frac e\varepsilon\right).
%\end{aligned}
%\]


For \(\lambda\in A_\eta\), this gives
\[
\begin{aligned}
        |h_\varepsilon(\lambda)|
        &\le
        |h_\varepsilon(1)|
        +
        |\lambda-1|
        \sup_{|\zeta-1|\le\eta}|h_\varepsilon'(\zeta)|        \\
        &\le
        C\varepsilon
        +
        C\eta\log \left(\frac{1}{\varepsilon}\right)
        \le
        C\varepsilon .
\end{aligned}
\]
Also,
\[
        |f(\lambda)-f(1)|
        \le C|\lambda-1|
        =
        C\eta
        =
        C\varepsilon^2,
        \qquad \lambda\in A_\eta .
\]
Thus
\[
        |E_\varepsilon(\lambda)|\le C\varepsilon,
        \qquad \lambda\in A_\eta .
\]
Since \(\operatorname{length}(A_\eta)\le 2\pi \eta\), we obtain
\[
\begin{aligned}
        \int_{A_\eta}
        |E_\varepsilon(\lambda)|
        \|(\lambda -T)^{-1}\|\,|d\lambda|
        &\le
        C\varepsilon
        \int_{A_\eta} \frac{|d\lambda|}{|\lambda-1|}        \\
        &\le
        C\varepsilon\, \frac{1}{\eta}\operatorname{length}(A_\eta)
        \le
        C\varepsilon .
\end{aligned}
\]
Hence the closing-arc contribution also tends to zero as \(\epsilon \downarrow 0\).

%Combining the estimates on \(B_\varepsilon\) and \(A_\eta\), 
Combining the estimates on the two parts of
\(\mathcal C_\varepsilon\), we get
%\[
%        h_\varepsilon(T)-f(T)+f(1)I
 %       \longrightarrow 0,
 %       \qquad \varepsilon\downarrow0 .
%\]
%Equivalently,
\[
        h_\varepsilon(T)\longrightarrow f(T)-f(1), \qquad  \varepsilon\downarrow0. 
\]
\end{proof}

\end{document}
